\begin{helper}[The terminal PE stored pair is disjoint from an earlier equal-level pair]
Let \(\mathcal E\) be a positive-rung maximum ETT on a finite loopless
multigraph, with displayed tree \(T\), certified series \(\mathcal T\), and
rung number \(n\).  Assume lower-rung assertions (i) and (ii) for every positive-rung
maximum ETT on the same input with fewer than \(n\) rungs: every proper
exactly-\(e_0\)-uncolored stable coloring of its final series tree makes any
two colors, one of which is missing on that tree, interchangeable there;
and its displayed vertex set is elementary under its generating coloring.

Let \(a=\operatorname{Fin.last}(n)\), and let \(c=r^+\ne a\) be an earlier
slot.  Suppose \(a\) and \(c\) are PE slots with distinct chosen supports
\(v_n,v_r\), both belonging to \(V(T)\), and suppose their dependent
\noderef{ETTFirstAppearanceIndex} values, computed using the supplied
membership witnesses, are equal.  Then their exact raw stored pairs are
disjoint:
\[
 \mathcal T.\mathrm{stored}(a)\cap
 \mathcal T.\mathrm{stored}(c)=\varnothing.
\]
\end{helper}
\begin{proof}
Write the earlier slot as \(h=c\).  Eliminating the PE certificates supplied
by \noderef{TashkinovSeries} and \noderef{TashkinovSeriesStep} identifies
the two raw stored sets as
\[
 S_h=\{\delta_h,\gamma_h\},\qquad
 S_n=\{\delta_n,\gamma_n\}.                                   \tag{1}
\]
We prove that both colors of \(S_h\) occur on internal edge identities at
\(v_n\) under \(\varphi_n\), and then exclude both terminal colors there.

In the second stable coloring \(\pi'_{h-1}\) used by the earlier PE, let
\(P\) and \(Q\) be the supports of \noderef{TwoColorComponent} through
\(v_h\) and \(v_n\), respectively, for the colors
\(\gamma_h,\delta_h\) in \noderef{TwoColorIncidenceGraph}.  The earlier PE
certificate puts \(v_h\) in \(T_h\), so
\noderef{ETTFirstAppearanceIndex} gives \(m(v_h)\le h\).  Equality of first
appearances and the ordered prefixes in \noderef{TashkinovSeriesStep} and
\noderef{TreeSequenceInitialSegment} therefore put \(v_n\) in \(T_h\).

The earlier \noderef{TreeExit} path, by \noderef{TreeExitPath}, runs from an
outside vertex to \(v_h\) and meets \(T_h\) only there.
\noderef{ParallelExtensionColor} makes \(\gamma_h\) missing at \(v_h\),
while \noderef{MaximumDefectiveChoice} supplies the incident
\(\delta_h\)-colored connecting edge and stability preserves it in
\(\pi'_{h-1}\).  Thus \(v_h\) has degree one in the two-color incidence
graph.  The outside endpoint likewise has degree one, because it misses one
of the colors and the path supplies the other.  Every internal path node
already has two neighbors, and
\noderef{TwoColorIncidenceDegreeAtMostTwo} forbids a third.  The shortest
attachment argument therefore shows that this exit path exhausts \(P\).
It meets \(T_h\) only at \(v_h\).  Since \(v_n\ne v_h\), \(P\ne Q\), and
\noderef{TwoColorComponentsDisjoint} makes them vertex-disjoint.
\noderef{TwoColorComponentSwap} changes only \(P\), so \(Q\) remains the
component through \(v_n\) under \(\varphi_h\).

The switch makes \(\delta_h\) missing at \(v_h\), hence by
\noderef{TreeMissingColors} it is missing on the final tree \(T_{h+1}\) of
the rung-\(h\) restriction.  That restriction is a
positive-rung maximum ETT because \(h\) is a positive PE slot, and
\noderef{CertifiedETTRestriction} says that all of its retained data agree
with the original series.  The assumed lower
assertion (ii), together with the definition
\noderef{TreeInterchangeable}, says that \(P\) is the unique
\((\gamma_h,\delta_h)\)-path component meeting \(T_{h+1}\).

We now show that \(Q\) is a cycle.  The earlier certificate says \(T_h\) is
\noderef{ColoredVertexSetClosed} under \(\varphi_{h-1}\), and that the
selected and second colorings are proper, exactly-\(e_0\)-uncolored, and
linked by the two displayed stability relations.  Applying
\noderef{sc1} in both directions transports any hypothetical boundary color
missing on \(T_h\) successively back to \(\varphi_{h-1}\), contradicting
closure.  Thus \(\pi'_{h-1}\) is closed on \(T_h\).

If \(h>1\), apply \noderef{CertifiedETTRestriction} at rung \(h-1\) and
then lower assertion (i); if \(h=1\), apply
\noderef{TechnicalTheoremBase} to the initial tuple.  In either case
\(T_h\) is \noderef{ColoredVertexSetElementary} under
\(\varphi_{h-1}\).  The two vertex clauses of
\noderef{StableColoring} preserve the individual missing sets, so \(T_h\)
is elementary under \(\pi'_{h-1}\).  Since \(\gamma_h\) is missing at
\(v_h\ne v_n\), it is not missing at \(v_n\).  By
\noderef{MissingColors} and \noderef{PresentColors}, an incident
\(\gamma_h\)-colored edge exists at \(v_n\).  Closedness and the exact
\(\delta_h\)-colored boundary connecting edge similarly show that
\(\delta_h\) is not missing at \(v_n\), so another incident edge has that
color.  The adjacency rule in \noderef{TwoColorIncidenceGraph} puts both
edge-nodes in \(Q\), and the disjoint switch component \(P\) does not alter
them.  Therefore \(Q\) is nonisolated and has both colored incidences at
\(v_n\).  The incidence-degree bound says it is a path or a cycle.  If it
were a path, it and \(P\) would be two distinct path components meeting
\(T_{h+1}\), contrary to interchangeability.  Hence \(Q\) is a cycle.

Let \(t>h\) be the first later slot that is not RE; it exists because the
terminal slot is PE.  The priority fields of
\noderef{RevisitingExtensionCertificate} in the certificate at \(t\) say that every
eligible two-color cycle from the preceding RE run is contained in the
ladder tree \(T_t\); hence \(Q\subseteq G[T_t]\).  Repeated
\noderef{HkuMissingStoredMonotone} gives
\[
 \overline\varphi_h(T_{h+1})\cup D_h
 \subseteq \overline\varphi_{t-1}(T_t)\cup D_{t-1}.             \tag{2}
\]
Both \(\gamma_h,\delta_h\) lie in the left side: the PE certificate puts
the raw stored pair inside its controlled union, and the switch makes
\(\delta_h\) missing on \(T_{h+1}\).  Now apply
\noderef{SameColorSeriesPreserved} to every exact edge of the internal cycle
\(Q\).  The internal-edge alternative excludes every connecting-edge
exception by \noderef{InternalEdges}, so all its edge colors are unchanged from \(\varphi_h\) through
\(\varphi_n\).  Consequently \(Q\) remains a
\((\gamma_h,\delta_h)\)-cycle through \(v_n\) under \(\varphi_n\).  At
\(v_n\) it supplies two distinct internal incident edges, one of each color.

Finally the terminal certificate and
\noderef{HkuParallelBoundaryStructure} identify the connecting edge at
\(v_n\) as the unique \(\gamma_n\)-colored boundary edge after the terminal
switch.  Directly from \noderef{ParallelExtensionColor},
\noderef{MaximumDefectiveChoice}, and
\noderef{TwoColorComponentSwap}, the switch replaces the unique incident
\(\delta_n\)-edge by a \(\gamma_n\)-edge and leaves no incident
\(\delta_n\)-edge; thus \(\delta_n\) is missing at \(v_n\).  Properness
prevents either internal cycle edge at \(v_n\) from having color
\(\gamma_n\), while missingness prevents either from having color
\(\delta_n\).  Hence
\(\{\gamma_h,\delta_h\}\cap\{\gamma_n,\delta_n\}=\varnothing\), which is
exactly the claimed disjointness by (1).
\end{proof}
